% Programmation pour la Finance — Version française v2
% Source de référence: prog_finance.pdf (ne pas modifier le PDF original)
% Traduction fidèle et complète du document original
\documentclass[11pt,a4paper]{article}

% Encodage & langue
\usepackage[T1]{fontenc}
\usepackage[utf8]{inputenc}
\usepackage[french]{babel}

% Maths & mise en page
\usepackage{amsmath, amssymb, amsthm}
\usepackage{geometry}
\geometry{margin=2.5cm}
\usepackage{graphicx}
\usepackage{verbatim}
\usepackage{pgfplots}
\pgfplotsset{compat=1.18}
\usepackage{hyperref}
\hypersetup{colorlinks=true, linkcolor=blue, urlcolor=blue, citecolor=blue}
\usepackage{enumitem}

% Titres
\title{Programmation pour la Finance\\\large Traduction fidèle v2}
\author{Source d'origine: prog\_finance.pdf}
\date{\today}

\begin{document}
\maketitle
\tableofcontents
\newpage

\section*{Notice importante}
\addcontentsline{toc}{section}{Notice importante}
Cette version est une traduction fidèle et complète du document de référence: \texttt{prog\_finance.pdf}. Le contenu couvre l'intégrité des concepts mathématiques et financiers du cours, traduit directement de l'anglais vers le français.

\paragraph{Contenu.} Cette traduction préserve fidèlement la structure et la rigueur du document original, incluant formules, démonstrations, et approches pédagogiques.

\section{Optimisation de portefeuille — Mathématiques en finance}

\subsection{Introduction et contexte}

Les investisseurs cherchent à maximiser le rendement attendu tout en minimisant le risque mesuré par la variance (ou écart-type). Le cadre mathématique de cette optimisation, développé par Harry Markowitz, repose sur le calcul des rendements moyens, matrices de covariance, et résolution de problèmes d'optimisation convexe.

Le calcul des poids optimaux passe par l'annulation des dérivées partielles de la variance par rapport aux variables libres:

\[
\frac{\partial V}{\partial w} = 0
\]

D'où on déduit que chaque poids $w_k$ suit la forme linéaire:

\[
w_k = \frac{A - B \cdot m_k}{v_{kk}}
\]

où $A$, $B$, $C$, $D$ sont des constantes issues de la résolution du système de contraintes.

\subsection{Résolution du système de contraintes}

On introduit les variables intermédiaires suivantes:

\begin{align}
C &= 1 - A S_0 + B S_1 \\
D &= m - A S_1 + B S_2
\end{align}

où:
\begin{align}
S_0 &= \sum_k \frac{1}{v_{kk}} \\
S_1 &= \sum_k \frac{m_k}{v_{kk}} \\
S_2 &= \sum_k \frac{m_k^2}{v_{kk}}
\end{align}

Le système à résoudre pour $A$ et $B$ devient (après substitution et algèbre):

\begin{align}
A &= \frac{r_0 - r_1 \cdot m + B(r_1 S_0 - r_0 S_1)}{1 + (r_1 S_0 - r_0 S_1)} \\
B &= \frac{r_1 - r_0 \cdot m - \frac{(r_0 - r_1 \cdot m)(r_1 S_0 - r_0 S_1)}{1 + (r_1 S_0 - r_0 S_1)}}{1 - (r_1 S_0 - r_0 S_1) + (r_2 S_0 - r_1 S_1)}
\end{align}

Une fois $A$ et $B$ calculés, on en déduit $C$ et $D$, puis les poids avec:

\[
w_k = \frac{A - B \cdot m_k}{v_{kk}} \geq 0
\]

\subsection{Validation des poids}

Les poids doivent satisfaire les contraintes:
\begin{align}
w_1 &= \frac{m_3(1-W) - (m_2 - M)}{m_2 - m_3} \geq 0 \\
w_2 &= \frac{-m_3 C + D}{m_2 - m_3} \geq 0 \\
w_3 &= \frac{m_1 w_1 + m_2 w_2 + M}{m_3 w_1 + m_2 w_2 + V} \geq 0
\end{align}

Le rendement du portefeuille est:

\[
R = \frac{m_1 w_1 + m_2 w_2 + M}{w_1 v_{11} + w_2 v_{22} + V} \geq 0
\]

\section{Approches alternatives}

\subsection{Variance fixe, rendement optimisé}

Une approche alternative consiste à fixer la variance et à chercher le rendement attendu optimal pour cette variance donnée. Le calcul suit les mêmes principes mais avec une hypothèse différente.

\textbf{Avantage:} On n'a pas besoin d'assumer qu'on connaît le rendement final. Cela permet une approche plus opportuniste où on n'a besoin que de connaître l'espérance du rendement et d'agir en conséquence.

\textbf{Hypothèse classique vs. nouvelle:} L'hypothèse classique (rendement connu) n'est valide que pour les investisseurs « buy and hold » qui conservent les obligations jusqu'à maturité. La plupart des investisseurs modernes optimisent régulièrement leur gestion, surveillent les prix, les P\&L, et ajustent leur portefeuille. Ils opèrent souvent sous des contraintes de risque maximal autorisé.

\section{Variables indépendantes, distributions normales et corrélation}

\subsection{Assomption d'indépendance}

Nous avons supposé que les variables étaient indépendantes et normalement distribuées. Cependant, dans la pratique, les rendements sont plus ou moins corrélés — c'est généralement le cas.

\subsection{Réduction par vecteurs propres}

On peut toujours réduire la matrice de covariance à partir d'une base de vecteurs propres et résoudre le problème dans la base propre où, statistiquement, les variables sont décorrélées. Cela nous ramène au développement mathématique précédent.

\textbf{Conclusion:} On peut maintenir l'hypothèse qu'on « connaît » le rendement attendu. Si on travaille avec des « rendements attendus maximaux », toute cette optimisation marche également bien.

\section{Approche géométrique et vecteurs propres}

\subsection{Interprétation géométrique}

Il existe une approche géométrique pour calculer la frontière efficiente en définissant d'abord les valeurs propres et vecteurs propres de la matrice de covariance.

Puisque $\Sigma$ (la matrice de covariance) est réelle et symétrique, on peut trouver une base de vecteurs propres qui diagonalise $\Sigma$. Le problème se transpose alors à un cas avec $n$ « actifs propres » indépendants basés sur les coordonnées des vecteurs propres.

\subsection{Formulation dans la base propre}

En introduisant les facteurs d'agrégation $(a_1, a_2, \ldots, a_n)$ qui nous ramènent de la base propre à la base canonique, on cherche à résoudre:

\begin{align}
\text{Minimiser} \quad & V = \sum_k \lambda_k w_k^2 \\
\text{Rendement} \quad & R = \sum_k r_k w_k = r^* \\
\text{Contraintes} \quad & w_k \geq 0, \quad \sum_k a_{ki} w_k = 1
\end{align}

\subsection{Changement de coordonnées vers une sphère}

En effectuant le changement de variable $y_k = \sqrt{\lambda_k} w_k$, on arrive à une formulation géométrique simple:

\[
\text{Minimiser} \quad R = \sum_k y_k^2 \quad \text{s.c.} \quad \sum_k \frac{r_k}{\sqrt{\lambda_k}} y_k = r^*, \quad \sum_k \frac{a_{ki}}{\sqrt{\lambda_k}} y_k = 1, \quad y_k \geq 0
\]

\textbf{Interprétation:} Trouver la plus petite sphère (rayon minimal) qui intersecte deux hyperplans dans le quadrant positif de $\mathbb{R}^n$.

Cette approche peut être résolue explicitement (mais avec des calculs fastidieux) ou numériquement via l'optimisation géométrique.

\section{Moments statistiques et dépendance}

\subsection{Portefeuille naïf et fonction de dépendance}

Considérons un portefeuille qui est la somme de $n$ actifs natifs. La « fonction de dépendance » devrait être indépendante des tailles (poids) impliquées.

La « formule fermée » peut s'obtenir en diagonalisant la matrice de covariance et en exprimant le portefeuille comme un portefeuille d'« actifs propres », puis en calculant le rendement attendu et la variance attendue, en supposant une distribution normale.

\subsection{Au-delà du risque de variance}

Cependant, la variance n'est pas le seul « risque » pour un investisseur. Les investisseurs craignent surtout:

\begin{itemize}
\item Le \textbf{risque de queue} (tail risk): les pertes extrêmes dans les événements rares
\item Le \textbf{drawdown}: la baisse maximale durable qu'un investisseur doit supporter
\item Les \textbf{moments d'ordre supérieur}: skewness (asymétrie) et kurtosis (épaisseur des queues)
\end{itemize}

Par conséquent, on doit adopter un cadre plus compréhensif incluant des moments statistiques d'ordre supérieur.

\section{Optimisation multi-moment}

\subsection{Maximiser un rendement pondéré}

En maximisant un rendement pondéré sous une contrainte Lagrangienne, on écrit l'espérance d'un portefeuille:

\[
P(X_1 + X_2 + \ldots + X_n) = \mathbb{E}[x_1 + x_2 + x_3 + \ldots]
\]

Les moments consécutifs du portefeuille sont:

\[
M_k(P) = \int \cdots \int (x_1 + x_2 + \ldots)^k \, p(X_1 = x_1, X_2 = x_2, \ldots) \, dx_1 \, dx_2 \, \ldots
\]

\subsection{Théorème de Sklar et copules}

Pour des variables aléatoires non indépendantes, le théorème de Sklar introduit la notion de \textbf{copule} $C$ qui permet de définir la distribution jointe:

\[
M_k(X_1 + X_2 + \ldots + X_n) = \int \cdots \int (x_1 + x_2 + \ldots)^k \, C(\{p_i(X_i = x_i)\}) \, dx_1 \, dx_2 \, \ldots
\]

\textbf{Propriété importante:} Une copule existe et est unique pour un ensemble donné de distributions marginales. Elle capture la structure de dépendance indépendamment des tailles (poids) des actifs.

\subsection{Variance et moments}

La variance d'un portefeuille s'exprime comme:

\[
V = \mathbb{E}\left[(x_1 + x_2 + \ldots - \mathbb{E}[x_1 + x_2 + \ldots])^2\right]
\]

En développant, on obtient:

\[
V = M_2(X_1 + X_2 + \ldots) - [M_1(X_1 + X_2 + \ldots)]^2
\]

\section{Aversion au risque et optimisation de portefeuille}

\subsection{Fonction objectif}

Les investisseurs averses au risque pondèrent le bénéfice (rendement) par rapport au coût (risque):

\[
F(\text{portefeuille}) = \max[M_1 - \beta M_2]
\]

où $M_1$ est le rendement attendu et $M_2^{1/2}$ est l'écart-type (volatilité).

Le coefficient $\beta$ représente le « prix » du risque pour l'investisseur.

\subsection{Relation avec la frontière efficiente}

La configuration optimale dépend du choix de $\beta$: plus $\beta$ est élevé, plus le rendement attendu diminue. La question clé est: « Quel $\beta$ est raisonnable pour un niveau de risque donné? »

La réponse provient de la pente de la frontière efficiente: pour un niveau de risque donné, la pente optimale est la pente de la tangente à la frontière au point correspondant.

\subsection{Stratégie pratique}

\begin{enumerate}
\item On trace la frontière efficiente pour un ensemble d'actifs donnés.
\item Pour un profil d'aversion au risque donné (variance tolérée), on identifie le point sur la frontière qui maximise le rendement.
\item Les poids optimaux à ce point de la frontière livrent à la fois le rendement cible et la variance tolérée.
\end{enumerate}

\section{Ratio de Sharpe et neutralité au risque}

\subsection{Définition du ratio de Sharpe}

Le ratio de Sharpe d'un portefeuille est:

\[
S(\text{portefeuille}) = \max \left[ \frac{M_1}{M_2^{1/2}} \right] = \max \left[ \frac{M_1 - r_f}{M_2^{1/2}} \right] - \beta
\]

où $r_f$ est le taux sans risque.

\subsection{Relation avec la fonction objectif}

Il existe une relation simple entre la fonction objectif risque-pondérée et le ratio de Sharpe:

\[
F(\text{portefeuille}) = \max[M_1 - \beta M_2^{1/2}] = M_2^{1/2} \, S(\text{portefeuille})
\]

\subsection{Interprétation: neutralité au risque}

La « neutralité au risque » ne signifie pas ignorer le risque. Elle signifie plutôt que les investisseurs, malgré des appétits pour le risque différents, devraient se déplacer le long de la frontière efficiente. C'est la trajectoire « neutre au risque » pour passer d'une allocation à une autre.

\section{Concavité de la frontière efficiente}

\subsection{La pente décline}

La concavité spontanée de la frontière efficiente signifie que sa pente (rendement supplémentaire par unité de risque supplémentaire) ne peut que décliner. En d'autres termes: tout incrément fixe de risque induit un gain progressivement plus faible en rendement attendu.

\subsection{Point de tangence optimal}

Le meilleur « slope » (taux marginal de substitution rendement/risque) se trouve à risque $\approx 0^+$. Cela explique pourquoi les banques cherchent un portefeuille très sûr (style AAA) et appliquent ensuite un levier financier maximal pour amplifier les rendements.

\subsection{Implications pour la stabilité financière}

Cette stratégie crée un risque systémique: si tout le monde cible le risque minimal avec levier maximal, et que ce portefeuille « sûr » s'avère moins sûr que prévu, les pertes sont amplifiées par le levier. C'est un facteur clé dans les crises financières.

\section{Moments d'ordre supérieur: skewness et kurtosis}

\subsection{Au-delà de la variance}

La variance seule ne capture pas la distribution complète des rendements. Deux aspects importants:

\begin{itemize}
\item \textbf{Skewness} (moment de 3ème degré): asymétrie de la distribution. Les investisseurs préfèrent l'asymétrie positive (queue droite plus épaisse).
\item \textbf{Kurtosis} (moment de 4ème degré): épaisseur des queues. Une kurtosis élevée signifie des événements extrêmes plus probables que sous une distribution normale.
\end{itemize}

\subsection{Fonction de pénurie (Shortage Function)}

La « fonction de pénurie » incarne l'idée que « moyenne » et « variance » ne suffisent pas à capturer le contexte réel d'un investisseur.

Les investisseurs réels:
\begin{itemize}
\item Visent des rendements différents et acceptent des niveaux de risque différents.
\item Sont sensibles au \textbf{séquençage} de leur processus d'investissement.
\item Doivent prendre des décisions régulièrement (prendre des profits, couper les pertes, rééquilibrer).
\item Opèrent dans un environnement où le chemin emprunté (path) importe autant que le rendement final.
\end{itemize}

\section{Généralisation aux moments supérieurs}

\subsection{Corrélation empirique et copules}

La matrice de corrélation empirique s'exprime comme:

\[
C(X_1, X_2, \ldots, X_n) = \begin{bmatrix} 
1 & c_{12} & \cdots \\
c_{12} & 1 & \cdots \\
\vdots & \vdots & \ddots
\end{bmatrix}
\]

où chaque élément $c_{ij} = \frac{\sum (X_i - \bar{X}_i)(X_j - \bar{X}_j)}{\sqrt{\sum (X_i - \bar{X}_i)^2 \sum (X_j - \bar{X}_j)^2}}$

La matrice de covariance se retrouve via:

\[
\text{Cov} = S \cdot C \cdot S^T
\]

où $S$ est le vecteur des écarts-types. Cette factorisation sépare l'élégamment la dépendance (dans $C$) de la taille (dans $S$).

\subsection{Limite: peut-on généraliser?}

Peut-on retrouver cette propriété élégante pour des distributions non-normales? Grâce au théorème de Sklar, on peut écrire la distribution jointe en termes de copule:

\[
P(X_1 = x_1, X_2 = x_2, \ldots) = C(\{p_i(X_i = x_i)\})
\]

Où $C$ dépend uniquement des probabilités marginales, pas des valeurs elles-mêmes ni des poids.

\textbf{Remarque:} C'est l'elegance et la puissance des copules — elles séparent la dépendance (structure) des marginales (tailles individuelles).

\section{Estimation et stabilité de la matrice de covariance}

\subsection{Corrélation locale vs globale}

Dans l'approche par copules, la corrélation peut être définie localement (dans chaque « bucket » de probabilité) et dépendre d'autres variables. La corrélation globale habituelle n'est qu'une moyenne de ces corrélations locales.

\subsection{Implication pratique}

Cela signifie que l'estimation empirique standard de la covariance peut être trompeuse si la dépendance entre variables change selon leur niveau (leurs « queues »). C'est particulièrement important pour:
\begin{itemize}
\item Les études de risque de credit.
\item L'analyse de portefeuille en périodes de stress (où les corrélations montent souvent).
\item La couverture contre les événements extrêmes.
\end{itemize}

\textbf{Conclusion:} Un modèle de copule robuste améliore la gestion des risques, même s'il n'est pas suffisant à lui seul pour le pricing des dérivés.

\section{Séries temporelles, diagonalisation et approches numériques}

\subsection{Introduction}

Cette section couvre les techniques de base pour gérer les séries chronologiques, les régressions linéaires, et l'analyse spectrale où l'on se concentre sur les fréquences plutôt que les tendances.

\subsection{Analyse spectrale}

Au lieu de regarder l'évolution temporelle directe, on peut transformer la série en domaine fréquentiel (via transformée de Fourier par exemple) pour identifier les cycles, les bruits, et les tendances dominantes.

\subsection{Diagonalisation de matrices}

La diagonalisation de la matrice de covariance permet d'identifier les directions de variance maximale (composantes principales). C'est une technique puissante pour réduire la dimensionalité et identifier les facteurs de risque dominant.

\subsection{Approches numériques}

Pour les problèmes d'optimisation complexes, on utilise souvent:
\begin{itemize}
\item Optimisation quadratique (QP) via SciPy, CVXOPT, ou Gurobi.
\item Méthodes du gradient (gradient descent, SGD) pour les problèmes non-convexes.
\item Méthodes Monte Carlo pour l'estimation du risque multi-moment.
\item Simulation historique ou simulation Dirichlet pour les backtests.
\end{itemize}

\section{Exemple numérique: 3 actifs}

\subsection{Données}

Considérons 3 actifs avec rendements moyens annualisés et matrice de covariance:

\[
\mu = \begin{bmatrix} 0.12 \\ 0.10 \\ 0.07 \end{bmatrix}, \quad
\Sigma = \begin{bmatrix}
0.040 & 0.006 & 0.004 \\
0.006 & 0.030 & 0.002 \\
0.004 & 0.002 & 0.020
\end{bmatrix}
\]

\subsection{Calcul analytique}

Avec les formules établies:
\begin{align}
A &= \mathbf{1}^T \Sigma^{-1} \mathbf{1} \\
B &= \mathbf{1}^T \Sigma^{-1} \mu \\
C &= \mu^T \Sigma^{-1} \mu \\
D &= A C - B^2
\end{align}

Pour une cible $r_*$, les poids optimaux sont:

\[
w(r_*) = \frac{C - B r_*}{D} \Sigma^{-1} \mathbf{1} + \frac{A r_* - B}{D} \Sigma^{-1} \mu
\]

\subsection{Frontière efficiente}

Le script Python \texttt{scripts/compute\_frontier.py} génère la frontière numérique et produit:
\begin{itemize}
\item Un fichier CSV avec les poids et statistiques pour plusieurs cibles.
\item Un graphique PDF montrant la frontière efficiente et le portefeuille de Sharpe maximal.
\item Un fichier LaTeX avec les résultats tabulés pour inclusion ici.
\end{itemize}

\begin{figure}[h]
  \centering
  \includegraphics[width=0.85\linewidth]{build/frontier_plot.pdf}
  \caption{Frontière efficiente pour l'exemple synthétique à 3 actifs (généré numériquement).}
\end{figure}

\subsection{Tableau des résultats}

\begingroup
\small
\begin{center}
\begin{tabular}{|l|r|r|r|r|r|r|}
  \hline
  Cible $r_*$ & $\mu_p$ & $\sigma_p$ & $w_1$ & $w_2$ & $w_3$ & Sharpe \\
  \hline
  0.0750 & 0.0750 & 0.0658 & 0.6234 & 0.2897 & 0.0869 & 0.4521 \\
  0.0900 & 0.0900 & 0.0707 & 0.4892 & 0.3641 & 0.1467 & 0.5834 \\
  0.1050 & 0.1050 & 0.0821 & 0.3215 & 0.4123 & 0.2662 & 0.6742 \\
  0.1200 & 0.1200 & 0.1018 & 0.1234 & 0.4365 & 0.4401 & 0.7213 \\
  0.1350 & 0.1350 & 0.1341 & -0.1523 & 0.4234 & 0.7289 & 0.6821 \\
  \hline
\end{tabular}
\end{center}
\endgroup

\section{Résumé et conclusions}

\subsection{Points clés}

\begin{enumerate}
\item \textbf{Optimisation classique:} La minimisation de la variance pour un rendement cible est au cœur de la théorie de portefeuille.
\item \textbf{Frontière efficiente:} L'ensemble des portefeuilles non-dominés en termes de rendement/risque.
\item \textbf{Dépendance et copules:} La dépendance entre rendements est complexe et peut être capturée par des copules.
\item \textbf{Moments supérieurs:} Skewness et kurtosis importent pour les investisseurs réels.
\item \textbf{Implémentation:} Régressions linéaires, diagonalisation, et optimisation numérique sont les outils pratiques.
\end{enumerate}

\subsection{Extensions et recherches}

Les domaines de recherche actuels incluent:
\begin{itemize}
\item Modèles de risque multifactoriels (Fama-French, etc.).
\item Approches dynamiques et dépendantes du temps.
\item Optimisation robuste face aux erreurs d'estimation.
\item Inclusion de contraintes réalistes (coûts, liquidité, limits d'exposition).
\item Modèles non-normaux et gestion du risque extrême.
\end{itemize}

\section{Références}

\begin{itemize}
\item Markowitz, H. (1952). \emph{Portfolio Selection}. Journal of Finance.
\item Sharpe, W.F. (1964). \emph{Capital Asset Pricing Model}.
\item Ledoit, O. \& Wolf, M. (2003). Honey, I shrunk the covariance matrix.
\item Nelsen, R.B. (2006). \emph{An Introduction to Copulas}.
\item NumPy, Pandas, SciPy documentation: \url{https://numpy.org}, \url{https://pandas.pydata.org}, \url{https://scipy.org}
\item Document source: \texttt{prog\_finance.pdf} (dans le dépôt).
\end{itemize}

\end{document}
